%% =========================================================
%% Chapter: Random sources and sampling
%% Source: September 24, 2026 polynomial-PEPS preprint, Theorem 5.2.
%% These are generic probability results, not the distributed compression theorem.
%% Gaussian source calculations are documented separately after their integration.
%% =========================================================

\chapter{Random Sources and Sampling}\label{ch:random_sources}

Finite product measures determine the second moments of independent source
slots. Diagonal coefficient moments then determine the second moments of
rectangular matrix sums. Summing estimates over nonempty correction subsets
gives an explicit sample count and a deterministic realization.
These calculations occur in the proof of
\cite[Theorem~5.2]{OpenAI2026PolynomialPEPS}.

All integrals are Bochner integrals, with value zero when the integrand is not
integrable. Empty products have value one, and scalar inversion uses
$0^{-1}=0$. For a complex rectangular matrix
$A$, write $\lVert A\rVert_2^2=\sum_{r,c}|A_{rc}|^2$ for its
Hilbert--Schmidt square.

\section{Independent source slots}

\begin{definition}[Products of slot coefficients]\label{def:random_tensor_coefficients}
    \lean{IndependentCovariance.tensorCoefficient}
    \leanok
    Let $S$ be a finite set. For each $s\in S$, let $\Omega_s$ be a measurable
    space and let $I_s$ and $L_s$ be two index sets. Given complex functions
    $c_s(a,b;\cdot)$ on $\Omega_s$, define the coefficient associated with
    multiindices $i\in\prod_s I_s$ and $l\in\prod_s L_s$ by
    \begin{align}
        z_{il}(\omega)
         & =\prod_{s\in S}c_s(i_s,l_s;\omega_s),
        \label{eq:random_tensor_coefficients}
    \end{align}
    where $\omega\in\prod_s\Omega_s$. The ket and bra index sets may have
    different cardinalities.
    This is the coefficient product in
    \cite[Equation~(5.11)]{OpenAI2026PolynomialPEPS}.
    % Source: 04-compression.tex:394--406, eq:compression-product-covariance.
\end{definition}

\begin{theorem}[Factorization of product moments]\label{thm:random_product_moments}
    \lean{IndependentCovariance.integral_product_mul_conj}
    \lean{IndependentCovariance.integrable_product_mul_conj}
    \lean{IndependentCovariance.integrable_tensorCoefficient_mul_conj}
    \leanok
    \uses{def:random_tensor_coefficients}
    Let each $\mu_s$ be a $\sigma$-finite measure on $\Omega_s$, and set
    $\mu=\bigotimes_{s\in S}\mu_s$. For complex functions $f_s,g_s$,
    \begin{align}
        \int\!\Bigl(\prod_s f_s(\omega_s)\Bigr)
        \overline{\Bigl(\prod_s g_s(\omega_s)\Bigr)}\,d\mu(\omega)
         & =\prod_s\int f_s(x)\overline{g_s(x)}\,d\mu_s(x).
        \label{eq:random_product_moments}
    \end{align}
    If each $f_s\overline{g_s}$ is integrable, the product integrand is
    integrable. In particular, pair-product integrability of the slot
    coefficients implies pair-product integrability of $z_{il}$.
\end{theorem}

\begin{proof}\leanok
    Conjugation commutes with the finite product, so the integrand is
    $\prod_s(f_s(\omega_s)\overline{g_s(\omega_s)})$.
    The finite-product integral formula gives
    (\ref{eq:random_product_moments}); its integrability statement applies to
    the same individual factors.
\end{proof}

\begin{theorem}[Diagonal moments of tensor coefficients]\label{thm:random_tensor_moments}
    \lean{IndependentCovariance.tensor_covariance_of_diagonal}
    \lean{IndependentCovariance.product_weighted_sqrt}
    \lean{IndependentCovariance.tensor_covariance_weighted}
    \lean{IndependentCovariance.sum_product_weights_eq_one}
    \leanok
    \uses{def:random_tensor_coefficients}
    On the product of $\sigma$-finite measures, suppose the individual slot
    moments satisfy
    \begin{align}
        \int c_s(a,b;x)\overline{c_s(a',b';x)}\,d\mu_s(x)
         & =\delta_{aa'}\delta_{bb'}\nu_s(a,b).
        \label{eq:random_slot_moments}
    \end{align}
    Then the full multiindices have diagonal moments:
    \begin{align}
        \int z_{il}\overline{z_{i'l'}}\,d\mu
         & =\delta_{ii'}\delta_{ll'}\prod_s\nu_s(i_s,l_s).
        \label{eq:random_tensor_moments}
    \end{align}
    In particular, on probability spaces, let $k$ be a positive integer and
    let $p_s,q_s$ be nonnegative weights. Put
    $p_i=\prod_s p_s(i_s)$ and $\widetilde p_l=\prod_s q_s(l_s)$. If
    $\nu_s(a,b)=k^{-1}\sqrt{p_s(a)q_s(b)}$, then
    \begin{align}
        \mathbb E[z_{il}\overline{z_{i'l'}}]
         & =k^{-|S|}\sqrt{p_i\widetilde p_l}\,
        \delta_{ii'}\delta_{ll'}.
        \label{eq:random_weighted_moments}
    \end{align}
    For finite $I_s$, the equalities $\sum_a p_s(a)=1$ imply
    $\sum_i\prod_s p_s(i_s)=1$; this normalization identity requires no sign
    condition on the individual weights.
    The weighted moment formula is
    \cite[Equation~(5.11)]{OpenAI2026PolynomialPEPS}.
    % Source: 04-compression.tex:394--406, eq:compression-product-covariance.
    % The general weighted declaration also handles k=0 with total inversion.
\end{theorem}

\begin{proof}\leanok
    \uses{thm:random_product_moments}
    Apply (\ref{eq:random_product_moments}) and substitute
    (\ref{eq:random_slot_moments}). A product of Kronecker deltas is one
    precisely when both complete multiindices agree, and is zero otherwise.
    The nonnegative-weight identity
    $\prod_s\sqrt{p_s(i_s)q_s(l_s)}=\sqrt{p_i\widetilde p_l}$
    gives (\ref{eq:random_weighted_moments}). Finally, finite distributivity
    gives $\sum_i\prod_s p_s(i_s)=\prod_s\sum_a p_s(a)=1$.
\end{proof}

\begin{theorem}[Gaussian coefficients on independent slots]
    \label{thm:random_gaussian_slots}
    \lean{QICLean.ComplexGaussian.pairedProductProbability}
    \lean{QICLean.ComplexGaussian.pairedProductProbability_nonneg}
    \lean{QICLean.ComplexGaussian.sum_pairedProductProbability_eq_one}
    \lean{QICLean.ComplexGaussian.independentDensityLaw}
    \lean{QICLean.ComplexGaussian.independentDensityLawIsProbabilityMeasure}
    \lean{QICLean.ComplexGaussian.independentDensityCoefficient}
    \lean{QICLean.ComplexGaussian.integrable_independentDensityCoefficient}
    \lean{QICLean.ComplexGaussian.integrable_independentDensityCoefficient_mul_conj}
    \lean{QICLean.ComplexGaussian.integral_independentDensityCoefficient_mul_conj}
    \lean{QICLean.ComplexGaussian.integral_independentDensityCoefficient_eq_zero}
    \leanok
    Let $S$ be finite. At each slot $s$, let $A_s,C_s$ be finite endpoint
    supports with nonnegative weights $\lambda_s,\mu_s$. Take independent
    standard circular complex Gaussian coordinates $g^{(j,s)}_{ac}$, with
    $\mathbb E|g^{(j,s)}_{ac}|^2=1$, and independent Gaussian factors across
    slots. For $k>0$, define
    \begin{align}
        c_s((a,b),(c,d))
               & =(\lambda_s(a)\lambda_s(b)\mu_s(c)\mu_s(d))^{1/4}
        \left(\frac1k\sum_{j=0}^{k-1}
        g^{(j,s)}_{ac}\overline{g^{(j,s)}_{bd}}
        -\delta_{ab}\delta_{cd}\right),
        \label{eq:random_gaussian_slot}                            \\
        z_{il} & =\prod_s c_s(i_s,l_s).
        \label{eq:random_gaussian_product}
    \end{align}
    The product law is a probability measure, and both $z_{il}$ and every
    conjugate pair product are integrable. Set
    $p_i=\prod_s\lambda_s(i_s^1)\lambda_s(i_s^2)$ and
    $\widetilde p_l=\prod_s\mu_s(l_s^1)\mu_s(l_s^2)$.
    Their covariance is (\ref{eq:random_weighted_moments}). If each endpoint
    weight list sums to one, both product weight lists also sum to one.
    For nonempty $S$, $\mathbb E z_{il}=0$.
    These are the actual independent coefficients in
    \cite[Equations~(5.7), (5.8), and~(5.11)]{OpenAI2026PolynomialPEPS}.
    % Source: 04-compression.tex:311--340,390--407, Gaussian covariance,
    % slot variance, and product covariance. Integrability also holds for k=0.
\end{theorem}

\begin{proof}\leanok
    \uses{thm:random_product_moments, thm:random_tensor_moments}
    The centered Gaussian pairing formula gives
    $\mathbb E[(g_\alpha\overline{g_\beta}-\delta_{\alpha\beta})
                \overline{(g_\gamma\overline{g_\delta}-\delta_{\gamma\delta})}]
        =\delta_{\alpha\gamma}\delta_{\beta\delta}$.
    Distinct sample coordinates remove the off-diagonal sample terms, so
    the slot covariance is $k^{-1}$ times its quarter-weight square.
    Theorem~\ref{thm:random_tensor_moments} gives the full covariance.
    The Gaussian coefficients belong to $L^2$, which gives the required
    product integrability. Finite-product factorization gives the zero
    mean for nonempty $S$, and finite distributivity gives normalization.
\end{proof}

\section{Rectangular trace-norm comparisons}

\begin{theorem}[Attained duality for the rectangular trace norm]
    \label{thm:random_trace_duality}
    \lean{Matrix.rectangularTraceNorm}
    \lean{Matrix.rectangularTraceNorm_square}
    \lean{Matrix.rectangularTraceNorm_eq_sum_sqrt_eigenvalues}
    \lean{Matrix.norm_trace_conjTranspose_mul_le_rectangularTraceNorm}
    \lean{Matrix.exists_contraction_trace_conjTranspose_mul_eq}
    \leanok
    For a complex rectangular matrix $A$, let $\lVert A\rVert_1$ be the sum
    of its singular values. Equivalently,
    $\lVert A\rVert_1=\sum_j\sqrt{\lambda_j(A^\dagger A)}$; for square
    matrices this is the trace norm. Every rectangular contraction $U$
    of the same size satisfies
    $|\tr(A^\dagger U)|\leq\lVert A\rVert_1$, and some such $U$ satisfies
    $\tr(A^\dagger U)=\lVert A\rVert_1$.
    This is the trace-norm duality used in
    \cite[Equation~(5.13)]{OpenAI2026PolynomialPEPS}.
    % Source: eq:compression-exterior-contraction, 04-compression.tex:454--470.
\end{theorem}

\begin{proof}\leanok
    Diagonalize $A^\dagger A$ in an orthonormal basis $v_j$.
    The trace is a sum of column inner products, each of absolute value at
    most $\lVert Av_j\rVert\lVert Uv_j\rVert
        \leq\sqrt{\lambda_j(A^\dagger A)}$.
    Equality is attained by sending $v_j$ to
    $Av_j/\sqrt{\lambda_j(A^\dagger A)}$ for positive eigenvalues, and to
    zero on the kernel. The resulting columns are orthonormal on the
    positive eigenspaces, so this map is a contraction.
\end{proof}

\begin{theorem}[Two-weight quarter-power comparisons]
    \label{thm:random_quarter_weights}
    \lean{Matrix.rectangularTraceNorm_halfWeighted_le}
    \lean{Matrix.frobeniusNormSq_rpow_quarter_mul_le_one}
    \leanok
    For nonnegative diagonal probability matrices $\tau,\widetilde\tau$
    and an arbitrary rectangular matrix $W$ of compatible size,
    \begin{align}
        \lVert\tau^{1/2}W\widetilde\tau^{1/2}\rVert_1
         & \leq\lVert\tau^{1/4}W\widetilde\tau^{1/4}\rVert_2.
        \label{eq:random_quarter_comparison}
    \end{align}
    More generally, let $R,\widetilde R$ be arbitrary positive
    semidefinite trace-one matrices and let $O$ be a rectangular
    contraction between their spaces. Then
    \begin{align}
        \lVert\widetilde R^{1/4}OR^{1/4}\rVert_2^2 & \leq1.
        \label{eq:random_dimension_free}
    \end{align}
    Zero eigenvalues and different input dimensions are allowed. These
    are the comparisons in
    \cite[Equations~(5.15) and~(5.17)]{OpenAI2026PolynomialPEPS}.
    % Source: eq:compression-quarter-powers, eq:compression-dimension-free,
    % 04-compression.tex:484--527. The latter extends the paper's diagonal case.
\end{theorem}

\begin{proof}\leanok
    \uses{thm:random_trace_duality}
    A contraction's squared column sums and squared row sums are at most
    one. Cauchy--Schwarz with the two probability weights therefore gives
    $\lVert\tau^{1/4}U\widetilde\tau^{1/4}\rVert_2\leq1$.
    Pull the half-weights into the trace pairing and apply
    Hilbert--Schmidt Cauchy--Schwarz against $U$; attained duality gives
    (\ref{eq:random_quarter_comparison}). For general positive matrices,
    diagonalize each in its own orthonormal basis and conjugate $O$ by
    the corresponding unitaries. Hilbert--Schmidt norm and the
    contraction bound are preserved, reducing
    (\ref{eq:random_dimension_free}) to the diagonal calculation.
\end{proof}

\begin{theorem}[Exterior contractions and discarded registers]
    \label{thm:random_exterior_contraction}
    \lean{Matrix.l2_opNorm_kronecker_one_le}
    \lean{Matrix.trace_partialTraceRight_conjTranspose_mul}
    \lean{Matrix.rectangularTraceNorm_mul_conjTranspose_le}
    \lean{Matrix.rectangularTraceNorm_partialTraceRight_le}
    \lean{Matrix.rectangularTraceNorm_partialTraceRight_mul_conjTranspose_le}
    \leanok
    Let $Z$ be an arbitrary rectangular matrix. Let $K$ and
    $\widetilde K$ be contractions from its respective output and input
    spaces into a common physical space tensored with a discard space.
    Then
    \begin{align}
        \lVert\Tr_{\mathrm{discard}}(KZ\widetilde K^\dagger)\rVert_1
         & \leq\lVert Z\rVert_1.
        \label{eq:random_exterior_contraction}
    \end{align}
    The same norm bound holds for left and right contractions alone,
    and partial trace alone cannot increase the trace norm.
    This is \cite[Equation~(5.13)]{OpenAI2026PolynomialPEPS}.
    % Source: eq:compression-exterior-contraction, 04-compression.tex:454--470.
\end{theorem}

\begin{proof}\leanok
    \uses{thm:random_trace_duality}
    Pull a physical test contraction $M$ back to
    $\widetilde K^\dagger(M\otimes\Id)K$. Its operator norm is at most one,
    since tensoring with the identity preserves that norm bound.
    Trace-norm duality gives (\ref{eq:random_exterior_contraction}).
    Taking only the corresponding part of this calculation gives the
    separate multiplication and partial-trace bounds.
\end{proof}

\section{Rectangular matrix second moments}

\begin{theorem}[Second moments from diagonal coefficient moments]
    \label{thm:random_matrix_moments}
    \lean{ProbabilityTheory.integrable_normSq_sum}
    \lean{ProbabilityTheory.integral_normSq_sum_eq_sum}
    \lean{ProbabilityTheory.integrable_frobenius_norm_sq_sum}
    \lean{ProbabilityTheory.integral_frobenius_norm_sq_sum_eq_sum}
    \lean{ProbabilityTheory.integral_frobenius_norm_sum_le_sqrt}
    \leanok
    Let $J$ be finite, let $\mu$ be a measure, and let $z_j$ be complex
    functions with every product $z_i\overline{z_j}$ integrable. Suppose
    $\int z_i\overline{z_j}\,d\mu=\delta_{ij}v_i$, where $v_i\in\R$.
    For complex numbers $a_j$ and arbitrary rectangular matrices $A_j$,
    \begin{align}
        \int\Bigl|\sum_j z_j a_j\Bigr|^2\,d\mu
         & =\sum_j v_j|a_j|^2,
        \label{eq:random_scalar_moment}     \\
        \int\Bigl\lVert\sum_j z_j A_j\Bigr\rVert_2^2\,d\mu
         & =\sum_j v_j\lVert A_j\rVert_2^2.
        \label{eq:random_matrix_moment}
    \end{align}
    Both squared integrands are integrable. If $\mu$ is a probability
    measure, then
    \begin{align}
        \mathbb E\Bigl\lVert\sum_j z_j A_j\Bigr\rVert_2
         & \leq\sqrt{\sum_j v_j\lVert A_j\rVert_2^2}.
        \label{eq:random_mean_norm}
    \end{align}
    These are the moment calculation and expectation estimate used in
    \cite[Equations~(5.16) and~(5.18)]{OpenAI2026PolynomialPEPS}.
    % Source: eq:compression-block-second-moment and eq:compression-one-choice,
    % 04-compression.tex:502--533. The weighted trace-norm comparison is separate.
\end{theorem}

\begin{proof}\leanok
    Expand $|\sum_j z_j a_j|^2$ into the finite sum
    $\sum_{i,j}z_i\overline{z_j}a_i\overline{a_j}$.
    Pair-product integrability permits integration term by term, and the
    diagonal moment identity removes every term with $i\neq j$.
    Apply (\ref{eq:random_scalar_moment}) to each matrix entry and sum over
    rows and columns to obtain (\ref{eq:random_matrix_moment}).
    On a probability space, the nonnegative variance of
    $f=\lVert\sum_j z_j A_j\rVert_2$ gives
    $(\mathbb E f)^2\leq\mathbb E(f^2)$, which implies
    (\ref{eq:random_mean_norm}).
\end{proof}

\section{Correction sums and deterministic realizations}

\begin{theorem}[Binomial correction cost]\label{thm:random_binomial_cost}
    \lean{CompressionSampling.sum_nonempty_subsets_pow}
    \lean{CompressionSampling.sum_nonempty_subsets_weight}
    \lean{CompressionSampling.inverse_half_power}
    \lean{CompressionSampling.sum_nonempty_subsets_rpow}
    \lean{CompressionSampling.zero_slots_error}
    \lean{CompressionSampling.binomial_error_le_exp}
    \lean{CompressionSampling.exp_eighth_sub_one_le_quarter}
    \leanok
    Let $T$ be a finite set of source positions, with $N=|T|$. For every real
    $a$,
    \begin{align}
        \sum_{\varnothing\neq U\subseteq T}a^{|U|}
         & =(1+a)^N-1.
        \label{eq:random_binomial_cost}
    \end{align}
    For $k>0$, the summand $Q^{|U|}k^{-|U|/2}$ equals
    $(Q/\sqrt{k})^{|U|}$, so its sum is
    $(1+Q/\sqrt{k})^N-1$. The sum is zero when $T$ is empty.
    If $a\geq0$, then $(1+a)^N-1\leq e^{Na}-1$.
    For $0\leq\epsilon\leq1$, one also has
    $e^{\epsilon/8}-1\leq\epsilon/4$.
    These are the scalar estimates in
    \cite[Equation~(5.19) and the following paragraph]{OpenAI2026PolynomialPEPS}.
    % Source: eq:compression-total-error, 04-compression.tex:540--565.
\end{theorem}

\begin{proof}\leanok
    Finite distributivity gives
    $\sum_{U\subseteq T}a^{|U|}=(1+a)^N$; remove the empty subset's contribution
    one. The square-root identity $k^{-n/2}=(\sqrt{k})^{-n}$ gives the
    sampling form. The inequality $1+a\leq e^a$ gives the exponential
    estimate. Finally, $|e^t-1|\leq2|t|$ for $|t|\leq1$, applied to
    $t=\epsilon/8$, gives the quarter-accuracy estimate.
\end{proof}

\begin{theorem}[A positive sample count with quarter accuracy]
    \label{thm:random_sample_count}
    \lean{CompressionSampling.sampleCount}
    \lean{CompressionSampling.sampleCount_bounds}
    \lean{CompressionSampling.ratio_le_of_sample_count}
    \lean{CompressionSampling.binomial_error_le_quarter}
    \lean{CompressionSampling.error_le_quarter_of_sample_count}
    \lean{CompressionSampling.error_le_quarter_sampleCount}
    \leanok
    For $N\in\N$ and real $Q,\epsilon$, put
    \begin{align}
        k & =\left\lceil(8NQ/\epsilon)^2\right\rceil+1.
        \label{eq:random_sample_count}
    \end{align}
    Then $k>0$ and
    $(8NQ/\epsilon)^2\leq k\leq(8NQ/\epsilon)^2+2$.
    If $Q\geq0$ and $0<\epsilon\leq1$, every positive integer $k$ satisfying
    $(8NQ/\epsilon)^2\leq k$ obeys
    \begin{align}
        (1+Q/\sqrt{k})^N-1 & \leq\epsilon/4.
        \label{eq:random_quarter_accuracy}
    \end{align}
    For $N>0$, the sample-count hypothesis gives
    $Q/\sqrt{k}\leq\epsilon/(8N)$.
    This is an explicit integer choice for the sample count in
    \cite[Theorem~5.2, paragraph following Equation~(5.19)]{OpenAI2026PolynomialPEPS}.
    % Source: 04-compression.tex:552--565. The ceiling-plus-one also covers Q=0.
\end{theorem}

\begin{proof}\leanok
    \uses{thm:random_binomial_cost}
    The ceiling inequalities give the three bounds on
    (\ref{eq:random_sample_count}). If $N=0$, the error is zero.
    Otherwise, the sample-count hypothesis gives
    $8NQ/\epsilon\leq\sqrt{k}$, hence $NQ/\sqrt{k}\leq\epsilon/8$.
    Apply the exponential and quarter-accuracy estimates of
    Theorem~\ref{thm:random_binomial_cost}.
\end{proof}

\begin{theorem}[Expected error of correction sums]\label{thm:random_expected_corrections}
    \lean{CompressionSampling.integral_error_le_binomial}
    \lean{CompressionSampling.integral_norm_sum_le_binomial}
    \lean{CompressionSampling.integral_norm_sum_le_sampling_error}
    \leanok
    Let $T$ be finite and let $F_U$ be integrable functions with values in a
    normed additive group, indexed by the nonempty subsets of $T$.
    If $\int\lVert F_U\rVert\,d\mu\leq a^{|U|}$ for each such $U$, then
    \begin{align}
        \int\Bigl\lVert\sum_{\varnothing\neq U\subseteq T}F_U\Bigr\rVert\,d\mu
         & \leq(1+a)^{|T|}-1.
        \label{eq:random_expected_corrections}
    \end{align}
    More generally, the same bound holds for an integrable real error $e$
    whenever $e\leq\sum_U e_U$ pointwise and the integrable real correction
    estimates satisfy $\int e_U\,d\mu\leq a^{|U|}$.
    In particular, for $k>0$, per-subset bounds
    $\int\lVert F_U\rVert\,d\mu\leq Q^{|U|}k^{-|U|/2}$ give the right side
    $(1+Q/\sqrt{k})^{|T|}-1$.
    This is the summation implication in
    \cite[Equation~(5.19)]{OpenAI2026PolynomialPEPS}.
    % Source: eq:compression-total-error, 04-compression.tex:540--551.
    % The per-subset estimates are explicit hypotheses, not discharged here.
\end{theorem}

\begin{proof}\leanok
    \uses{thm:random_binomial_cost}
    Apply the triangle inequality pointwise, integrate the finite sum, and
    substitute the individual integral bounds. Formula
    (\ref{eq:random_binomial_cost}) evaluates their sum. For a real error,
    integral monotonicity replaces the triangle inequality.
\end{proof}

\begin{theorem}[A deterministic realization with half accuracy]
    \label{thm:random_realization}
    \lean{CompressionSampling.exists_error_le_half_sampleCount}
    \lean{CompressionSampling.exists_valid_error_le_half_sampleCount}
    \lean{CompressionSampling.exists_valid_norm_sum_le_half_sampleCount}
    \leanok
    Let $\mu$ be a probability measure, let $Q\geq0$ and
    $0<\epsilon\leq1$, and choose $k$ by (\ref{eq:random_sample_count}).
    Suppose $e$ is integrable and
    $\mathbb E e\leq(1+Q/\sqrt{k})^N-1$.
    If a condition $V$ holds almost surely, there exists $\omega$ satisfying
    $V(\omega)$ and $e(\omega)\leq\epsilon/2$.
    In particular, this conclusion applies with
    $e=\lVert\sum_{\varnothing\neq U\subseteq T}F_U\rVert$ and $N=|T|$,
    provided each $F_U$ is integrable and
    $\mathbb E\lVert F_U\rVert\leq Q^{|U|}k^{-|U|/2}$.
    This is the realization step after
    \cite[Equation~(5.19)]{OpenAI2026PolynomialPEPS}.
    % Source: 04-compression.tex:558--567. No distributed-network conclusion is tagged.
\end{theorem}

\begin{proof}\leanok
    \uses{thm:random_sample_count, thm:random_expected_corrections}
    There is a point outside the exceptional null set of $V$ at which
    $e\leq\mathbb E e$. Theorem~\ref{thm:random_sample_count} bounds this
    expectation by $\epsilon/4\leq\epsilon/2$.
    For correction sums, Theorem~\ref{thm:random_expected_corrections}
    supplies the required expectation bound.
\end{proof}
